\begin{abstract}
\noindent Whether transcriptional responses to unmeasured genetic perturbations can be predicted better than by simple baselines is contested, and the answer depends heavily on evaluation design. We built an evaluation protocol for Perturb-seq data with four elements: batch-matched responses; disjoint control cells for training and evaluation targets; technical reproducibility limits from disjoint batches; and tests of how well each metric separates a technical duplicate from uninformative predictions. We developed models on one screen, then fixed an internal analysis plan before applying them to further screens. Four met the pre-specified inclusion criteria: two cell lines from one large CRISPRi study, and CRISPRa and CRISPRi screens in iPSC-derived neurons. We evaluated simple model classes (linear and nearest-neighbour baselines, and small neural-network ensembles using screen-derived co-expression, target baseline expression or Gene2vec features); graph neural networks and single-cell foundation models were not tested. For genes absent from training, these models outperformed the mean of other perturbations only in the two large screens, reducing its squared error by about 5\%. This held across output gene sets, control-cell splits, initialisation seeds and a bootstrap over co-expressed gene modules. Within those screens, gains increased with the number of training perturbations. Between a quarter and two-thirds of perturbations in the large screens had specific responses that did not reproduce across technical batches. The ability of mean squared error to separate informative from uninformative predictions tracked this reliability closely, so low discrimination reflects measurement noise rather than a defect of the metric. Rankings of held-out genes by predicted response were enriched for strong responders 1.5- to 3-fold in the large screens; whether this would improve experimental decisions remains untested. With the information and model classes examined, unseen-perturbation prediction in these screens is weak, concentrated in reproducible perturbations, and easy to mask or overstate through the choice of metric and baseline.
\end{abstract}
